\chapter{Requirements Analysis}
\label{ch:requirements}

\section{Project Objectives and Constraints}

The requirements are expressed at system level so that a model metric is not
mistaken for a complete security property. \Cref{tab:reqs} records the evidence
currently available. ``Partial'' means that an artefact exists but the full
requirement has not been demonstrated.

{\small
\begin{longtable}{@{}p{13mm}p{66mm}p{25mm}p{41mm}@{}}
\caption{Requirement and evidence matrix.}\label{tab:reqs}\\
\toprule
ID & Requirement & Status & Current evidence or gap \\
\midrule
\endfirsthead
\toprule
ID & Requirement & Status & Current evidence or gap \\
\midrule
\endhead
FR-01 & Capture a short facial video and validate basic media quality before a
decision. & Partial & Flutter clip capture exists; a production quality gate is
not established. \\
FR-02 & Run a lightweight temporal first-stage model locally without sending all
media to a server. & \implemented & MobileNetV3 PyTorch Lite model and Flutter
inference flow use a 24-frame contract. \\
FR-03 & Produce an interpretable risk score and decision state rather than only a
class label. & Partial & A probability and fixed 0.5 decision are exposed; score
calibration and policy bands are absent. \\
FR-04 & Escalate uncertain or policy-defined high-risk sessions to stronger
analysis. & \proposed & Candidate server models exist independently; routing is
not integrated into the application. \\
FR-05 & Detect both physical presentation and digital manipulation attacks. &
Partial & Both families are studied, but no unified operational test set exists. \\
FR-06 & Detect camera bypass or digitally injected media. & \proposed & Threat is
documented; sensor attestation/IAD is not implemented. \\
FR-07 & Route a frame to one of five manipulation specialists using embedding
evidence. & Partial & Six-way legacy router exists; intended no-original router
requires revised labels and rerun. \\
FR-08 & Send unresolved cases to accountable human review with reason codes and
appeal. & \proposed & Workflow and interface are not implemented or evaluated. \\
NFR-01 & Preserve subject/video-group separation in evaluation and prevent
training leakage. & Partial & FF++ routing split is group-disjoint; every dataset
protocol still requires a documented audit. \\
NFR-02 & Measure FAR, FRR, HTER, AUC, calibration, and subgroup behaviour under
deployment-like conditions. & Partial & HTER/AUC and routing metrics exist;
calibration and demographic analysis are absent. \\
NFR-03 & Protect biometric data in transit, at rest, and throughout retention and
deletion. & \proposed & Required control; no production data plane exists. \\
NFR-04 & Bound latency, memory, battery, bandwidth, and server cost. & Partial &
Mobile integration and a small G2V2 latency pilot exist; full device profiling is
absent. \\
NFR-05 & Log model/version/decision metadata without retaining unnecessary raw
biometrics. & \proposed & Audit schema and retention schedule are not implemented. \\
NFR-06 & Remain usable under poor capture, low bandwidth, accessibility needs,
and legitimate appearance variation. & \proposed & Needs user and field testing. \\
\bottomrule
\end{longtable}
}

Key constraints are the lack of private deployment data, limited GPU access,
heterogeneous dataset licences and annotations, mobile compute limits, and the
changing attack landscape. The most important safety constraint is epistemic: a
score produced on a public benchmark cannot be represented as fraud probability
for a real applicant without deployment-domain validation and calibration.

\section{Regulatory Requirements and Standards}

\begin{table}[H]
\centering
\caption{Standards and guidance relevant to the target system.}
\label{tab:standards}
\begin{tabularx}{\textwidth}{@{}p{33mm}Y Y@{}}
\toprule
Source & Relevance & Correct use in this report \\
\midrule
ISO/IEC 30107-3:2023 \cite{iso30107} & PAD testing and reporting, including
attack-presentation types. & Evaluation reference; no certification or
conformance is claimed. Injection outside the capture device is out of scope. \\
NIST SP 800-63A-4 \cite{nist80063a4} & Remote proofing, PAD, genuine-sensor
confidence, protected channels, injection and forged-media controls. & Security
design guidance; it is not presented as binding Bangladeshi law. \\
ENISA remote-ID reports \cite{enisa2022,enisa2024} & Presentation/injection threat
models and layered countermeasures. & Threat and architecture guidance, not a
product certification. \\
ACM/IEEE ethical codes \cite{acmcode,ieeecode} & Avoid harm and discrimination;
respect privacy, honesty, competence, and accountability. & Basis for professional
responsibility and reporting limitations. \\
Bangladesh legal and financial requirements & Biometric processing, cyber
security, consumer rights, financial records, and appeal. & Must be reviewed with
qualified legal/compliance personnel before deployment; no legal conclusion is
made in this draft. \\
\bottomrule
\end{tabularx}
\end{table}

The design should implement purpose limitation, data minimisation, informed and
accessible notice, least-privilege access, encryption, deletion enforcement,
incident response, and auditable override/appeal. These controls cannot be added
only at the end because model training itself processes sensitive face media.

\section{Formulation of Solutions}

Four solution patterns were considered. A mobile-only model offers low latency
and privacy but limited capacity and visibility into injection. A server-only
model permits stronger analysis but increases network dependence, data transfer,
cost, and central privacy risk. Active challenges add session-specific evidence
but create friction and can exclude some users. A risk-based cascade combines a
lightweight local path with selective escalation and capture controls.

\begin{table}[H]
\centering
\small
\setlength{\tabcolsep}{4pt}
\caption{Qualitative design decision matrix (5 is favourable). Scores are design
judgements and must be validated by field evidence.}
\label{tab:decision}
\begin{tabular}{@{}lccccc@{}}
\toprule
Alternative & Security breadth & Mobile latency & Privacy & Accessibility & Cost \\
\midrule
Mobile-only passive model & 2 & 5 & 4 & 4 & 5 \\
Server-only strong model & 3 & 2 & 2 & 3 & 2 \\
Mandatory active challenge & 4 & 3 & 3 & 2 & 3 \\
Risk-based layered cascade & 5 & 4 & 4 & 4 & 3 \\
\bottomrule
\end{tabular}
\end{table}

The selected target is the layered cascade in \Cref{fig:pipeline}. It keeps the
integrated mobile temporal model as the first analytical layer, proposes capture
integrity and quality checks before it, and escalates only according to calibrated
policy. Candidate second-layer evidence includes $G^2V^2$former, GD-FAS, a
universal Xception model, or the specialist router. Their outputs are not
interchangeable: they should be calibrated, tested for complementary error, and
fused only after a common protocol. Human review is reserved for unresolved cases
and is a governed fallback, not proof that automation is safe.

\begin{figure}[H]
  \centering
  \resizebox{\textwidth}{!}{\begin{tikzpicture}[
  font=\footnotesize,
  node distance=7mm and 8mm,
  box/.style={draw=buetblue,rounded corners,fill=softblue,align=center,minimum height=12mm,text width=25mm},
  impl/.style={box,fill=softgreen},
  prop/.style={box,fill=softorange},
  decision/.style={diamond,aspect=2,draw=buetblue,fill=white,align=center,inner sep=1.5pt,text width=18mm},
  arr/.style={-{Latex[length=2mm]},thick,draw=buetblue}
]
\node[prop] (integrity) {Capture integrity\\and quality gate};
\node[impl,right=of integrity] (mobile) {Layer 1: mobile\\temporal MobileNetV3\\24 frames};
\node[decision,right=of mobile] (conf) {risk /\\confidence};
\node[prop,above right=9mm and 16mm of conf] (pass) {Accept or reject\\by calibrated policy};
\node[prop,below right=9mm and 16mm of conf] (layer2) {Layer 2: stronger\\server analysis};
\node[prop,right=of layer2] (fusion) {Fusion and\\calibration};
\node[prop,right=of fusion] (review) {Human review\\for unresolved cases};

\draw[arr] (integrity) -- (mobile);
\draw[arr] (mobile) -- (conf);
\draw[arr] (conf) -- node[above,sloped]{clear} (pass);
\draw[arr] (conf) -- node[below,pos=0.43,align=center]{uncertain/\\high risk} (layer2);
\draw[arr] (layer2) -- (fusion);
\draw[arr] (fusion) -- (review);

\node[draw=green!50!black,rounded corners,fit=(mobile),inner sep=2mm,label={[green!40!black]below:integrated prototype}] {};
\node[draw=orange!80!black,dashed,rounded corners,fit=(integrity)(pass)(layer2)(fusion)(review),inner sep=3mm,
 label={[orange!70!black]above:target architecture}] {};
\end{tikzpicture}
}
  \caption{Layered system boundary. Green identifies the integrated prototype;
  orange identifies target-design elements that are not yet end-to-end.}
  \label{fig:pipeline}
\end{figure}
